% Supplemental Material for "Shape versus circuit: why graphene geometric
% diodes and Tesla valves do not make terahertz rectennas".
% Draft 2026-11-10. Numbers from results/*.md and results/checks/.
\documentclass[aps,prapplied,onecolumn,superscriptaddress,floatfix]{revtex4-2}
\usepackage{graphicx}
\usepackage{amsmath,amssymb}
\usepackage{array}
\usepackage[colorlinks=true,allcolors=blue]{hyperref}
\renewcommand{\thetable}{S\arabic{table}}
\renewcommand{\thesection}{S\arabic{section}}
\newcommand{\vF}{v_{F}}
\newcommand{\um}{\,\mu\mathrm{m}}
\newcommand{\nm}{\,\mathrm{nm}}

\begin{document}
\title{Supplemental Material: Shape versus circuit: why graphene geometric diodes and Tesla valves do not make terahertz rectennas}
\author{Todd Chapman}
\date{\today}
\maketitle

This supplement lists the acceptance tests that each solver passed before it
was used, the closing convergence checks, and the numerical parameters of the
production runs. Every test is a program in the public repository and prints
the measured number against its target; the test names below are those of the
code, with the plasmonic tests prefixed P to distinguish them from the
ballistic ones.

\section{Ballistic solver}

\begin{table}[ht]
\caption{Acceptance tests of the trajectory solver at dc (200\,000 particles
per run). All pass.}
\begin{ruledtabular}
\begin{tabular}{lp{0.52\textwidth}p{0.36\textwidth}}
Test & What is checked & Target \\
\hline
T1 & specular straight channel, no field & transmission exactly 1 \\
T2 & diffuse-wall channel, length to width 0.5 to 16 & two-dimensional free-molecular transmission factors~\cite{richley1964numerical} 0.804, 0.683, 0.541, 0.400, 0.274, 0.179 \\
T3 & asymmetric funnel, no field, $p=1$ and $0.3$, $\lambda=700\nm$ & $W_{S}T_{SD}=W_{D}T_{DS}$ within 3.5 standard errors \\
T4 & the same funnel with the field on, both injection sides at one total energy & $N_{S}T_{SD}=N_{D}T_{DS}$ \\
T5 & mirror-symmetric constriction under bias & asymmetry 1 \\
T6 & zero bias & zero injection weight and zero current \\
T7 & three-stage side-loop valve, diffuse walls & trapped, leaked and Dirac-cut fraction below $10^{-3}$ \\
T8 & ballistic $1\um$ channel at $10^{12}$ cm$^{-2}$ & Sharvin conductance 8.74 mS within 1 \% \\
T9 & $10\um$ channel, $\lambda=300\nm$ & $T=\ell_{1}/(L+\ell_{1})$, $\ell_{1}=\pi\lambda/2$, within 5 \% \\
T10 & chord tolerance 0.05 against 0.005 nm & same transmission \\
T11 & funnel of $100/10/40\nm$ (lead, neck, drain) at zero field & $T=0.0995$ against the neck-to-lead ratio 0.100 (information) \\
\end{tabular}
\end{ruledtabular}
\end{table}

\begin{table}[ht]
\caption{Acceptance tests of the time-dependent solver and the two-sided
estimator ($10^{6}$ particles per side; T16 $4\times10^{6}$, T17
$2\times10^{6}$). Currents in $\mu$A for the clean specular funnel at
$0.1$~V. All pass.}
\begin{ruledtabular}
\begin{tabular}{lp{0.47\textwidth}p{0.41\textwidth}}
Test & What is checked & Result \\
\hline
T12 & integration of $\hbar\dot{\bm k}=\bm F$ without energy re-imposition, static field, against the dc solver & $T(+0.1)=0.14141$ (reference $0.14130$), $T(-0.1)=0.04148$, $A=1.622$ (reference $1.624$) \\
T13 & two-sided equilibrium estimator against the bias-window estimator; zero bias & $188.7\pm0.6$ against $189.2\pm0.5$; $-0.15\pm0.66$ at $V=0$ \\
T14 & quasi-static anchor at 300 MHz: rectified current and $|Y_{1}|$ against the static curve & $17.78\pm0.59$ against $18.27\pm0.18$; $1.589\pm0.013$ against $1.580$ mS \\
T15 & straight control at 3 GHz and 1 THz & $-0.80\pm0.57$ and $-0.20\pm0.62$; trapped plus leaked below $2\times10^{-6}$ \\
T16 & chord cap $\vF/40f$ against $\vF/160f$, 500 nm drop, 1 THz & $6.32\pm0.30$ against $5.55\pm0.32$ \\
T17 & Dirac cut against a classical turning point, 500 nm drop, 1 THz & $5.94\pm0.45$ against $5.69\pm0.45$ \\
\end{tabular}
\end{ruledtabular}
\end{table}

The quasi-static anchor is at 300 MHz and not at a few gigahertz because of
the slow-drift population described in the main text. The equilibrium
injection of the two-sided estimator samples the kinetic distribution down
to zero energy, and between $10^{-3}$ and $7\times10^{-3}$ of the injected
electrons are driven below the 2 meV cut that stands in for the Dirac
point; these would return to their own contact and carry no current. In the
harmonic estimator a drain-injected electron removed by the cut must be
booked as returned, otherwise $|Y_{1}|$ at the anchor is $6\%$ low. Five
operating points recomputed with a classical turning point in place of the
cut agree with the production values within one standard error.

\section{Viscous solver}

\begin{table}[ht]
\caption{Acceptance tests of the steady Navier-Stokes solver. All pass.}
\begin{ruledtabular}
\begin{tabular}{lp{0.44\textwidth}p{0.44\textwidth}}
Test & What is checked & Result \\
\hline
H1 & Poiseuille flow, no slip and Navier slip $S=0.5$, 2, $\infty$ & exact to $10^{-12}$ \\
H2 & Brinkman-Poiseuille flow, $\mathrm{Da}=10$, $S=0.5$; Brinkman plug flow & velocity error $O(h^{2.7})$, $1.2\times10^{-6}$ at $h=25\nm$; plug flow exact \\
H3 & Stokes reciprocity of a side-loop valve, three $(\mathrm{Da},S)$ & $D-1=2\times10^{-12}$ to $2\times10^{-11}$; wall flux exactly zero \\
H4 & Kovasznay flow~\cite{kovasznay1948laminar}, $\mathrm{Re}=40$ & velocity error $O(h^{2.9})$, $2.2\times10^{-5}$ at $h=25\nm$ \\
H5 & backward-facing step~\cite{armaly1983experimental,erturk2008numerical}, expansion ratio 2 & reattachment at 2.90 step heights ($\mathrm{Re}=100$) and 4.96 ($\mathrm{Re}=200$) \\
H6 & side-loop valve as a classical no-slip device & $D=0.9993$, 1.0037, 1.0197 at $\mathrm{Re}=25$, 50, 100 \\
H7 & mesh convergence, side-loop valve, $\mathrm{Re}=10$ & $D-1=1.01$, 1.10, $1.14\times10^{-2}$ at wall spacing 20, 14, 10 nm \\
H8 & mirror-symmetric constriction, $\mathrm{Re}=10$ & $D-1=7\times10^{-6}$ \\
H9 & $(D-1)/\mathrm{Re}$ at $\mathrm{Re}=0.05$ to 0.2; mirrored and re-meshed device & constant to 2 \%; directional solver bias $6\times10^{-7}$ \\
U1 & units & Poiseuille and Brinkman resistances equal the Gurzhi and Drude formulas to $10^{-12}$ \\
\end{tabular}
\end{ruledtabular}
\end{table}

Below $\mathrm{Re}\approx1$ the diodicity is under the discretisation floor of
about $3\times10^{-5}$ and is reported as a bound. For the classical valve
with no-slip walls the steady branch continued from Stokes flow ends at a
fold near $\mathrm{Re}=85$; this is outside the range relevant to graphene.

\begin{table}[ht]
\caption{Mesh convergence of the classical Tesla valve ($175\nm$ channel, no
damping): $D-1$ at wall spacings of 20 nm (16\,703 nodes, the production
mesh), 14 nm (25\,447 nodes) and 10 nm (38\,449 nodes).}
\label{tab:mesh}
\begin{ruledtabular}
\begin{tabular}{llccc}
Wall condition & $\mathrm{Re}$ & 20 nm & 14 nm & 10 nm \\
\hline
perfect slip & 3 & 0.0119 & 0.0118 & 0.0124 \\
perfect slip & 10 & 0.234 & 0.232 & 0.243 \\
perfect slip & 30 & 1.472 & 1.510 & 1.619 \\
slip length $0.5\,W$ & 3 & 0.0068 & 0.0068 & 0.0072 \\
slip length $0.5\,W$ & 10 & 0.151 & 0.151 & 0.159 \\
slip length $0.5\,W$ & 30 & 1.145 & 1.177 & 1.257 \\
\end{tabular}
\end{ruledtabular}
\end{table}

Table~\ref{tab:mesh} shows that the diodicity of the classical valve
changes by less than $3\%$ between the two coarser meshes and by $5$ to
$7\%$ between the two finer ones; the pressure drop itself converges only
at first order because of the re-entrant corners at the loop junctions. The
values quoted in the main text are from the 20 nm mesh and are good to
about $10\%$.

\section{Plasmonic solver}

\begin{table}[ht]
\caption{Acceptance tests of the perturbation solver (P1 to P9, P15 to P17)
and of the time-domain solver (P10 to P14). All pass.}
\begin{ruledtabular}
\begin{tabular}{lp{0.42\textwidth}p{0.46\textwidth}}
Test & What is checked & Result \\
\hline
P1 & static limit against the viscous solver & relative difference $5\times10^{-4}$ \\
P2 & linear plasmon response of a channel against $1/\cos kL$ & errors $7\times10^{-9}$ to $1.2\times10^{-5}$ \\
P3 & position and width of the first plasmon resonance, with friction, viscosity and viscoelastic viscosity & position 0.1568 (exact 0.1571); widths within 1 \% \\
P4 & photoresponse against the Dyakonov-Shur closed form~\cite{dyakonov1996detection} and an independent quadrature & errors $2\times10^{-7}$ to $2\times10^{-4}$ \\
P5 & funnel, quasi-static rectified current: curvature of the steady solution against perturbation & $2.206\times10^{-2}$ against $2.210\times10^{-2}$ \\
P6 & time domain against perturbation at small amplitude & ratios 0.9997 and 0.9994 \\
P7 & symmetric termination: dc current of a symmetric constriction and of a channel relative to the funnel & $6\times10^{-7}$ and $10^{-8}$ against $1.6\times10^{-3}$ \\
P8 & mesh convergence of the funnel response on three meshes & last change 2.0 \% and 1.1 \% at reduced frequency 1 and 3 \\
P9 & units & plasma speed, gate overdrive, resistance and frequency units against hand values \\
P10 & time-domain efficiency against the time-discrete perturbation law & ratios 1.0033 and 1.0038 \\
P11 & dc load correction from a deliberately wrong bias seed & 1.0016 (short circuit), 1.0017 (open circuit) \\
P12 & time step: 60 against 120 steps per period at $100$~K & efficiency ratio 0.991 \\
P13 & power bookkeeping & absorbed power within 0.2 \% of the linear value \\
P14 & P10 for the symmetric termination & ratio 0.998 \\
P15 & contact resistance: $I_{\mathrm{sc}}(R_{\mathrm{dc}}+2R_{c})=V_{\mathrm{oc}}$; port relations; passivity & $7\times10^{-15}$; $4\times10^{-16}$; absorbed fraction 0.39 \\
P16 & invariance under co-scaling of density, gate distance, width and scattering times at fixed nondimensional groups & nondimensional response identical to round-off \\
P17 & units and closure at a second operating point & closed forms to $2\times10^{-16}$; Dirac against shallow-water response ratio 0.970 \\
\end{tabular}
\end{ruledtabular}
\end{table}

The time-domain runs start from the perturbative state of the
time-discretised operator, which removes a start-up transient that would
otherwise exceed the rectified signal at small amplitude; the dc load
current is held at the small-signal prediction and corrected with the dc
Thevenin relation. The level of the square law carries the 60-step
time-discretisation bias ($+0.6\%$, $+4.5\%$ and $-0.5\%$ at 0.1, 0.3 and
1 THz) and the $2\%$ mesh uncertainty of P8; the ratios to the square law
quoted in the main text are taken against the time-discrete law, so the
bias cancels. The shape of the response is converged in the time step: at
$100$~K the efficiency ratio between 60 and 120 steps per period is 0.991
(P12), and at $300$~K the ratio of the efficiency to the square law changes
from 0.992, 1.054, 1.101 to 1.004, 1.058, 1.109 at amplitudes 0.1, 0.5, 0.7
of the voltage unit at 0.3 THz, and from 1.144 to 1.145 at 0.7 at 0.1 THz.

\section{Production parameters}

Ballistic dc: 400\,000 trajectories per bias point in 64 batches. Ballistic
ac: $10^{6}$ trajectories per contact and point ($4\times10^{6}$ at
$\lambda=700\nm$ and at the smaller amplitude), chords limited to $40\nm$, to
$\vF/40f$ and to one twentieth of the field region. Viscous: Taylor-Hood
meshes with $20\nm$ wall spacing ($23\,000$ to $55\,000$ triangles),
Newton residual below $10^{-10}$. Plasmonic small signal: wall spacing
$14\nm$ for the $175\nm$ devices and $80\nm$ ($18\,601$ nodes) for the
$1\um$ funnel, 25 frequencies from 10 GHz to 3 THz. Plasmonic large signal:
the same $1\um$ mesh, 60 steps per period, eight amplitudes from $0.01$ to
$0.7$ of the voltage unit.

\bibliography{refs,refs_detectors}

\end{document}
